\question{4.1}{
    Een kioskhouder verkoopt elke zaterdag een aantal exemplaren van het \emph{Weekendmagazine}.
    Het aantal exemplaren dat op een willekeurige zaterdag wordt gevraagd, kan worden beschouwd als een kansvariabele $X$ waarvoor de volgende uitkomsten en bijbehorende kansen gegeven zijn (zie de tabel).

    \begin{center}
        \begin{tabular}{cc}
            \toprule
                {\bfseries Aantal gevraagde \emph{Weekendmagazines} $k$} & {\bfseries $P(X=k)$} \\
            \cmidrule{1-1} \cmidrule{2-2} 
                $30$ & $0.10$ \\
                $35$ & $0.35$ \\
                $40$ & $0.30$ \\    
                $45$ & $0.20$ \\
                $50$ & $0.05$ \\
            \bottomrule
        \end{tabular}
}
\begin{enumerate}[label=(\alph*)]
    \item De kioskhouder koopt elke zaterdag $40$ exemplaren van het \emph{Weekendmagazine}.
    Hoe groot is de kans dat hij niet aan de vraag kan voldoen?
    \answer{
        De kioskhouder kan niet aan de vraag voldoen als er meer dan $40$ exemplaren worden gevraagd.
        Dit is het geval (zie tabel) als er $45$ of meer kranten worden gevraagd.
        De gevraagde kans luidt:
        \[
            P(X = 45) + P(X = 50) = 0.20 + 0.05 = 0.25.
        \]
    }

    \item Hoe groot is de verwachtingswaarde van de vraag op een willekeurige zaterdag?
    \answer{
        Volgens de definitie van verwachtingswaarde geldt de volgende formule:
        \begin{align*}
            E(X)    &= \sum_k k \cdot P(X=k).
        \end{align*}
        Term voor term invullen levert:
        \[
            E(X) = 30 \cdot 0.10 + 35 \cdot 0.35 + 40 \cdot 0.30 + 45 \cdot 0.20 + 50 \cdot 0.05 = 38.75.        
        \]
    }

    \item Hoe groot is de variantie van de vraag?
    \answer{
        De formule voor de variantie luidt:
        \[
            \Var{X} = \sum_k (k - E(X))^2 \cdot P(X = k).
        \]
        Hiertoe berekenen we eerst de termen $(k - E(X))^2$:
        \begin{center}
            \begin{tabular}{cc}
                \toprule
                    \makecell{{\bfseries Aantal gevraagde} \\ {\bfseries \emph{Weekendmagazines} $k$}} & {\bfseries $(k - E(X))^2$} \\
                \cmidrule{1-1} \cmidrule{2-2} 
                    $30$ & $(30 - 38.75)^2 = 76.5625$ \\
                    $35$ & $(35 - 38.75)^2 = 14.0625$ \\
                    $40$ & $(40 - 38.75)^2 = 1.5625$ \\    
                    $45$ & $(45 - 38.75)^2 = 39.0625$ \\
                    $50$ & $(50 - 38.75)^2 = 126.5625$ \\
                \bottomrule
            \end{tabular}
        \end{center}
        Term voor term invullen levert:
        \begin{align*}
            \Var{X} &= \sum_k (k - E(X))^2 \cdot P(X = k) \\
                    &= 76.5625 \cdot 0.10 + 14.0625 \cdot 0.35 + \ldots + 126.5625 \cdot 0.05 \\
                    &= 21.1875. 
        \end{align*}
    }

    \item De inkoopprijs van de kranten is voor de kioskhouder \euro 1.40, terwijl de verkoopprijs \euro 2.25 bedraagt.
    De kranten die niet worden verkocht, moeten worden weggegooid (en leveren dus een verlies op van \euro 1.40 per stuk).
    Bereken de verwachte winst van de kioskhouder bij een inkoop van 40 exemplaren per zaterdag.
    \answer{
        We berekenen de opbrengst van de verschillende aanallen die gevraagd kunnen worden.
        De (onzekere) opbrengst per zaterdag definiëren we als de kansvariabele $Y$.
        Bij een vraag van $30$ stuks bedraagt de opbrengst $30 \cdot 0.85 + 10 \cdot (-1.40) = \text{\euro} 11.50$.
        Bij een vraag van $35$ stuks bedraagt de opbrengst $35 \cdot 0.85 + 5 \cdot (-1.40) = \text{\euro} 22.75$.  
        Bij een vraag van $40$ (of meer!) stuks bedraagt de opbrengst $40 \cdot 0.85 = \text{\euro} 34$.

        We hebben hier dus te maken met de volgende kansvariabele $Y$:
        \begin{center}
            \begin{tabular}{cc}
                \toprule
                    {\bfseries Winst $y$} & {\bfseries $P(Y = y)$} \\
                \cmidrule{1-1} \cmidrule{2-2} 
                    $11.50$ & $0.10$ \\
                    $22.75$ & $0.35$ \\
                    $34$ & $0.30 + 0.20 + 0.05 = 0.55$ \\    
                \bottomrule
            \end{tabular}
        \end{center}
        De verwachte winst is te berekenen met de formule voor de verwachtingswaarde:
        \begin{align*}
            E(Y)    &= \sum_y y \cdot P(Y = y) \\
                    &= 11.50 \cdot 0.10 + 22.75 \cdot 0.35 + 34 \cdot 0.55 \\
                    &\approx \text{\euro} 27.81.
        \end{align*}
    }
\end{enumerate}